% 2.8 — brief en documento/investigacion/2-8-identidad-y-acceso.md
\subsection{GESTIÓN DE IDENTIDAD Y ACCESO}
\subsubsection{OAUTH 2.0 Y OPENID CONNECT}

La gestión de identidad y acceso reúne dos procesos distintos: la autenticación, que identifica y valida a un usuario, y la autorización, que le concede acceso \parencite{KeycloakServerAdmin}. El estándar de referencia para el segundo es OAuth 2.0 (\textit{Open Authorization}), un marco que permite a una aplicación obtener acceso limitado a un servicio HTTP, en nombre del propietario de un recurso o por cuenta propia \parencite{RFC6749}. La especificación define cuatro roles: el propietario del recurso (\textit{resource owner}), que cuando es una persona se llama usuario final; el servidor de recursos (\textit{resource server}), que aloja los recursos protegidos y responde a peticiones que presentan tokens de acceso; el cliente, que hace esas peticiones en nombre del propietario; y el servidor de autorización (\textit{authorization server}), que emite los tokens tras autenticar al propietario. Un único servidor de autorización puede emitir tokens aceptados por múltiples servidores de recursos \parencite{RFC6749}, propiedad en la que se apoya el inicio de sesión único (\textit{single sign-on}, SSO) que adopta este proyecto.

OAuth 2.0 distingue clientes confidenciales, que mantienen en secreto sus credenciales porque se ejecutan en un servidor, y clientes públicos, como una aplicación nativa o una aplicación en navegador \parencite{RFC6749}. Sobre esa distinción se organizan sus cuatro concesiones (\textit{grants}). En el código de autorización (\textit{authorization code}) el cliente redirige al usuario al servidor de autorización, que lo autentica y devuelve un código que el cliente canjea por el token; las credenciales del usuario nunca se comparten con el cliente y el token no pasa por el navegador. En el flujo implícito el token se emite directamente al navegador, el servidor no autentica al cliente y el token puede quedar expuesto a otras aplicaciones. En la concesión por credenciales del propietario (\textit{resource owner password credentials}, ROPC) el cliente recibe directamente usuario y contraseña, y la de credenciales de cliente (\textit{client credentials}) se usa cuando el cliente actúa por cuenta propia y está reservada a clientes confidenciales \parencite{RFC6749}.

En un cliente público el código de autorización puede ser interceptado en un canal no protegido por TLS\@. La prueba de clave para el intercambio de código (\textit{Proof Key for Code Exchange}, PKCE) mitiga ese ataque: el cliente genera un verificador secreto, envía un derivado en la petición de autorización y presenta el verificador original al canjear el código, de modo que el servidor comprueba que quien canjea es quien inició el flujo \parencite{RFC7636}. La práctica actual de seguridad para OAuth 2.0 \parencite{RFC9700} lo vuelve obligatorio para los clientes públicos, incluidas las aplicaciones web; desaconseja el flujo implícito; establece que ROPC no debe usarse, porque expone las credenciales del propietario al cliente; y exige que cada servidor de recursos verifique en cada petición que el token fue emitido para él. El borrador de OAuth 2.1 \parencite{OAuth21Draft}, aún en elaboración en el IETF, consolida esa dirección al incorporar PKCE al código de autorización y omitir los flujos implícito y ROPC\@. \textcite{Fett2016} aportan el respaldo académico: en el primer análisis formal extensivo de OAuth 2.0 descubren cuatro ataques hasta entonces desconocidos y demuestran que, con las correcciones propuestas, el protocolo cumple autorización, autenticación e integridad de sesión. La seguridad del marco depende, por tanto, de implementar el conjunto completo de mitigaciones, lo que favorece adoptar un proveedor que ya las incorpora.

El resultado de estos flujos es un token de acceso, una cadena que representa una autorización y que suele ser opaca para el cliente, acompañada opcionalmente de un token de refresco para renovarlo \parencite{RFC6749}. El tipo más usado es el token portador (\textit{bearer}): cualquier parte que lo posea puede usarlo sin demostrar posesión de material criptográfico, por lo que debe protegerse en almacenamiento y en tránsito y enviarse en la cabecera \texttt{Authorization} \parencite{RFC6750}. Cuando el token es autocontenido, el formato habitual es el token web JSON (\textit{JSON Web Token}, JWT): una representación compacta de afirmaciones (\textit{claims}) sobre un sujeto, firmada digitalmente, con nombres registrados como emisor (\texttt{iss}), sujeto (\texttt{sub}), audiencia (\texttt{aud}) y expiración (\texttt{exp}) \parencite{RFC7519}. OAuth 2.0, sin embargo, no dice quién es el usuario: sin un perfilado adicional es incapaz de proporcionar información sobre la autenticación del usuario final \parencite{OIDCCore2014}. OpenID Connect (OIDC) es la capa de identidad construida sobre OAuth 2.0 que cubre ese vacío: el cliente, llamado parte confiante (\textit{relying party}), solicita el ámbito \texttt{openid} a un proveedor de OpenID (OP) y recibe, además del token de acceso, un token de identidad (\textit{ID token}) en formato JWT con afirmaciones sobre la autenticación realizada, que debe validar comprobando emisor y audiencia \parencite{OIDCCore2014}.

El cuadro~\ref{tab:oauth-proyecto} traslada estos roles y flujos al proyecto. El ERP de ISI Mustang autentica hoy con un endpoint propio que recibe usuario y contraseña y emite un JWT generado por el mismo servidor. Ese esquema sirve para una sola aplicación, pero el proyecto agrega un servicio de predicción en FastAPI, un servidor MCP y una API para Power BI, y \textcite{Almeida2022} muestran que en arquitecturas de varios servicios cada uno debe asegurarse individualmente, que el acceso no autorizado a uno puede comprometer el sistema completo, y que OAuth 2.0, OpenID Connect y JWT son los mecanismos documentados para ese escenario. Con un servidor de autorización externo ningún servicio recibe contraseñas, cada uno valida los tokens por firma y audiencia, y la identidad es la misma en todos. Se prevé que Angular opere como cliente público con código de autorización y PKCE, que las llamadas entre servicios usen credenciales de cliente, y que NestJS, FastAPI y el servidor MCP sean servidores de recursos con audiencias distintas.

\begin{table}[H]
  \centering
  \caption{Roles y flujos de OAuth 2.0 / OpenID Connect y su uso en el proyecto}\label{tab:oauth-proyecto}
  \begin{tabular}{@{}>{\RaggedRight\arraybackslash}p{0.22\textwidth}>{\RaggedRight\arraybackslash}p{0.34\textwidth}>{\RaggedRight\arraybackslash}p{0.38\textwidth}@{}}
    \toprule
    \textbf{Rol o flujo} & \textbf{Definición} & \textbf{En el proyecto} \\
    \midrule
    \multicolumn{3}{@{}l}{\textit{Roles}} \\
    Servidor de autorización / OP & Autentica al usuario y emite tokens de acceso e identidad & Keycloak (véase 2.8.2) \\
    Propietario del recurso & Usuario final que concede el acceso & Personal de ISI Mustang que usa el ERP \\
    Cliente público & Aplicación en navegador o nativa, sin secreto & Frontend Angular del ERP \\
    Cliente confidencial & Aplicación en servidor que guarda sus credenciales & NestJS al invocar a FastAPI; Power BI al consumir la API de NestJS; cliente del servidor MCP \\
    Servidor de recursos & Acepta tokens y sirve recursos protegidos & API NestJS, servicio FastAPI y servidor MCP, cada uno con audiencia propia \\
    \midrule
    \multicolumn{3}{@{}l}{\textit{Flujos}} \\
    Código de autorización + PKCE & Inicio de sesión interactivo de personas & Angular ante Keycloak \\
    Credenciales de cliente & Acceso máquina a máquina, solo clientes confidenciales & NestJS hacia FastAPI; Power BI hacia la API de NestJS \\
    Implícito & Token emitido directamente al navegador; desaconsejado & No se usa \\
    ROPC & El cliente recibe usuario y contraseña; no debe usarse & No se usa; sustituye al inicio de sesión propio del ERP \\
    \bottomrule
  \end{tabular}
  \fuente{Elaboración propia a partir de \textcite{RFC6749}, \textcite{RFC9700} y \textcite{OIDCCore2014}.}
\end{table}

\subsubsection{KEYCLOAK}

Keycloak es una solución de inicio de sesión único para aplicaciones web y servicios web RESTful cuyo objetivo es simplificar la seguridad: las funciones que los desarrolladores tendrían que escribir por sí mismos vienen incorporadas \parencite{KeycloakServerAdmin}. Es un servidor separado que se administra en la red propia; las aplicaciones se configuran para apuntar a él y ser aseguradas por él mediante protocolos abiertos como OpenID Connect o SAML 2.0. Es un proyecto de código abierto en incubación en la Cloud Native Computing Foundation, y \textcite{Thorgersen2023} le dedican una monografía que complementa la documentación oficial. En los términos de la sección anterior, Keycloak es a la vez servidor de autorización de OAuth 2.0 y proveedor de OpenID Connect: expone los endpoints de autorización, token, información de usuario, certificados e introspección, que las aplicaciones descubren a partir de un documento de configuración publicado por cada dominio \parencite{KeycloakSecuringApps}.

Su modelo se organiza en pocos conceptos \parencite{KeycloakServerAdmin}. Un dominio (\textit{realm}) gestiona un conjunto de usuarios, credenciales, roles y grupos; los usuarios inician sesión en un dominio y los dominios están aislados entre sí. Los clientes son las entidades que pueden pedir a Keycloak que autentique a un usuario: aplicaciones que quieren asegurarse y ofrecer inicio de sesión único, o servicios que solo necesitan un token de acceso para invocar a otros servicios protegidos. Los roles identifican una categoría de usuario (administrador, gerente, empleado) y las aplicaciones asignan permisos a roles antes que a usuarios individuales; la asignación de roles puede encapsularse en los tokens, de modo que cada aplicación decide los permisos sobre sus recursos sin consultar al servidor. Keycloak es además un proveedor de identidad (\textit{identity provider}, IdP) capaz de federar la autenticación hacia otros proveedores OIDC o SAML y hacia directorios LDAP, y soporta cualquier plataforma que cuente con una librería de parte confiante de OpenID Connect; \textcite{Thorgersen2023} confirman que los adaptadores propios de versiones anteriores fueron reemplazados por librerías estándar.

El aspecto más relevante para una arquitectura de varios servicios es cómo se validan los tokens. Los tokens de acceso de Keycloak son JWT firmados con JSON Web Signature, por lo que pueden validarse localmente con la clave pública del dominio emisor, obtenida del endpoint de certificados y almacenada en caché; consultar el endpoint de introspección en cada petición, en cambio, puede ser lento y sobrecargar el servidor \parencite{KeycloakSecuringApps}. Es el patrón que \textcite{Thorgersen2023} describen en su aplicación de ejemplo: el token de identidad identifica al usuario, el token de acceso transporta los permisos para invocar el backend, y este valida los tokens y controla el acceso por los roles que lee de ellos. Keycloak implementa el código de autorización, con el que la aplicación obtiene token de acceso, de refresco y de identidad, y las credenciales de cliente, que no forman parte de OpenID Connect sino de OAuth 2.0; lista los flujos implícito y ROPC como desaconsejados y recomienda registrar las URI de redirección lo más específicas posible, en especial para clientes públicos \parencite{KeycloakSecuringApps}.

Estas propiedades explican por qué el proyecto reemplaza la autenticación propia del ERP por Keycloak. Con el esquema actual la identidad vive dentro del backend en NestJS: cualquier servicio nuevo que quisiera reconocer al mismo usuario tendría que compartir el secreto con el que se firman los tokens o delegar la validación en el ERP, con el acoplamiento y el riesgo de compromiso en cadena que \textcite{Almeida2022} señalan. Con Keycloak, un único dominio contiene a los usuarios y roles de ISI Mustang; Angular autentica a las personas con código de autorización y PKCE; NestJS, FastAPI y el servidor MCP validan localmente el mismo JWT con la clave pública del dominio, comprueban su audiencia y autorizan por los roles que trae; y las llamadas entre servicios, y desde Power BI hacia la API de NestJS, usan credenciales de cliente. Se prevé desplegar Keycloak en contenedor junto al resto de la plataforma, en el servidor propio de la empresa; las librerías con las que cada servicio valida los tokens y la configuración del dominio son decisiones de implementación que se documentan en el Capítulo IV\@.
